\documentclass[aspectratio=169]{beamer}

\usepackage[utf8]{inputenc}
\usepackage[T1]{fontenc}
\usepackage[spanish]{babel}
\usepackage{amsmath,amssymb}
\usepackage{array}
\usepackage{booktabs}
\usepackage{tabularx}
\usepackage{ragged2e}
\usepackage{xcolor}

\usetheme{Madrid}
\usecolortheme{default}
\setbeamertemplate{navigation symbols}{}
\setbeamertemplate{blocks}[rounded][shadow=false]
\setlength{\columnsep}{0.7cm}

\definecolor{VerdeRiego}{RGB}{41,92,64}
\definecolor{ArenaRiego}{RGB}{173,120,46}
\definecolor{GrisSuave}{RGB}{245,247,244}

\setbeamercolor{structure}{fg=VerdeRiego}
\setbeamercolor{title}{fg=VerdeRiego}
\setbeamercolor{frametitle}{fg=VerdeRiego,bg=GrisSuave}
\setbeamercolor{block title}{fg=white,bg=VerdeRiego}
\setbeamercolor{block body}{bg=GrisSuave}
\setbeamercolor{alerted text}{fg=ArenaRiego}
\setbeamerfont{block title}{size=\normalsize,series=\bfseries}
\setbeamerfont{block body}{size=\small}

\title[Plan de riego]{Plan de Riego \'Optimo de una Finca}
\subtitle{Presentaci\'on resumida del informe}
\author[Equipo ADA II]{Juan Carlos Cruz Mu\~noz \and David Alejandro Enciso Gutierrez \and Juan Esteban Rodriguez Valencia \and Estiven Andres Martinez Granados}
\institute[Univalle]{Universidad del Valle \\ An\'alisis de Algoritmos II}
\date{Mayo de 2026}

\begin{document}

\begin{frame}
  \vspace{1.0cm}
  \begin{beamercolorbox}[wd=\textwidth,rounded=true,sep=0.6cm,center]{block title}
    {\color{white}\rmfamily\bfseries\Large Plan de Riego \'Optimo de una Finca}
  \end{beamercolorbox}

  \vspace{0.9cm}
  \begin{center}
    {\large \insertauthor\par}
    \vspace{0.8cm}
    {\normalsize \insertinstitute\par}
    \vspace{0.8cm}
    {\large \insertdate\par}
  \end{center}
\end{frame}

\begin{frame}{Problema y modelo}
  \begin{columns}[T,onlytextwidth]
    \column{0.40\textwidth}
    \small
    \begin{block}{Datos del tabl\'on}
      \begin{tabular}{@{}ll@{}}
        \toprule
        S\'imbolo & Significado \\
        \midrule
        \(ts_i\) & supervivencia \\
        \(tr_i\) & tiempo de riego \\
        \(p_i\) & prioridad \\
        \(rp_i\) & d\'ia perfecto \\
        \bottomrule
      \end{tabular}
    \end{block}

    \begin{alertblock}{Objetivo}
      \[
      \Pi^\star=\arg\min_{\Pi} CR_F^\Pi
      \]
    \end{alertblock}

    \column{0.56\textwidth}
    \scriptsize
    \begin{block}{Costo por tabl\'on al iniciar \(t_i^\Pi\)}
      \[
      c_i^\Pi=
      \begin{cases}
      ts_i-(t_i^\Pi+tr_i), & t_i^\Pi=rp_i,\\
      2\left(ts_i-(t_i^\Pi+tr_i)\right), &
      t_i^\Pi\le ts_i-tr_i,\ t_i^\Pi\neq rp_i,\\
      2p_i\big((t_i^\Pi+tr_i)-ts_i\big), &
      t_i^\Pi>ts_i-tr_i.
      \end{cases}
      \]
      \[
      CR_F^\Pi=\sum_{i=0}^{n-1} c_i^\Pi
      \]
    \end{block}

    \begin{block}{Idea central}
      El costo cambia de r\'egimen seg\'un dos referencias:
      \[
      rp_i
      \qquad\text{y}\qquad
      ts_i-tr_i.
      \]
    \end{block}
  \end{columns}
\end{frame}

\begin{frame}{Algoritmos comparados}
  \small
  \begin{table}
    \centering
    \begin{tabularx}{\textwidth}{@{}>{\RaggedRight\arraybackslash}p{2.7cm}>{\Centering\arraybackslash}p{2.6cm}>{\Centering\arraybackslash}p{1.8cm}X@{}}
      \toprule
      Enfoque & Complejidad & Exacto & Lectura r\'apida \\
      \midrule
      Fuerza bruta & \(O(n!\,n)\) & S\'i & Revisa todas las permutaciones; sirve como referencia exacta para \(n\) peque\~no. \\
      Voraz & \(\Theta(n\log n)\) & No & Muy r\'apido; construye una soluci\'on localmente razonable. \\
      Programaci\'on din\'amica & \(O(n2^n)\) & S\'i & Exacta y mucho m\'as escalable que fuerza bruta. \\
      \bottomrule
    \end{tabularx}
  \end{table}

  \vspace{0.4em}
  \begin{block}{Pregunta del proyecto}
    Comparar calidad de la soluci\'on, complejidad y correctitud de los tres enfoques.
  \end{block}
\end{frame}

\begin{frame}{Fuerza bruta: idea y papel pr\'actico}
  \begin{columns}[T,onlytextwidth]
    \column{0.47\textwidth}
    \scriptsize
    \begin{block}{Idea del algoritmo}
      \begin{enumerate}
        \item Generar todas las permutaciones de los \(n\) tablones.
        \item Evaluar el costo total de cada orden.
        \item Conservar la mejor permutaci\'on encontrada.
      \end{enumerate}
      En la implementaci\'on se usa \texttt{next\_permutation}.
    \end{block}

    \begin{alertblock}{Garant\'ia}
      Recorre todo el espacio factible y, por construcci\'on, devuelve una permutaci\'on de costo \'optimo.
    \end{alertblock}

    \column{0.47\textwidth}
    \scriptsize
    \begin{block}{Aplicaci\'on a los ejemplos}
      \begin{tabular}{@{}lccc@{}}
        \toprule
        Finca & \(n\) & Permutaciones & Costo \'optimo \\
        \midrule
        \(F_1\) & 5 & \(5!=120\) & 20 \\
        \(F_2\) & 5 & \(5!=120\) & 32 \\
        \bottomrule
      \end{tabular}
      \[
      \Pi^\star_{F_1}=\Pi^\star_{F_2}=\langle 2,1,0,3,4\rangle
      \]
    \end{block}

    \begin{block}{L\'imite}
      \[
      T_{\texttt{roFB}}(n)=O(n!\,n)
      \]
      Exacta, pero pr\'actica solo en instancias peque\~nas. En la comparaci\'on del proyecto se report\'o solo hasta \(n\le 8\).
    \end{block}
  \end{columns}
\end{frame}

\begin{frame}{Programaci\'on din\'amica: idea central}
  \begin{columns}[T,onlytextwidth]
    \column{0.48\textwidth}
    \small
    \begin{block}{Estado del subproblema}
      Implementaci\'on \textit{top-down} con memoizaci\'on. El estado queda determinado por el conjunto \(M\) de tablones ya regados.
      \[
      \tau(M)=\sum_{i\in M} tr_i
      \]
      \[
      DP(M)=\text{costo m\'inimo restante desde } M
      \]
    \end{block}

    \begin{block}{Recurrencia}
      \[
      DP(M)=
      \min_{i\notin M}
      \left\{
      c_i(\tau(M))+DP(M\cup\{i\})
      \right\}
      \]
      con caso base \(DP(U)=0\).
    \end{block}

    \column{0.46\textwidth}
    \small
    \begin{block}{Reconstrucci\'on}
      Se guarda en \(\texttt{eleccion}[M]\) el tabl\'on que minimiza la recurrencia y luego se reconstruye la permutaci\'on \'optima siguiendo:
      \[
      \varnothing \to M_1 \to M_2 \to \cdots \to U.
      \]
    \end{block}

    \begin{alertblock}{Complejidad}
      \[
      T_{\texttt{roPD}}(n)=O(n2^n),
      \qquad
      S_{\texttt{roPD}}(n)=O(2^n)
      \]
      Mucho mejor que fuerza bruta, aunque sigue siendo exponencial.
    \end{alertblock}
  \end{columns}
\end{frame}

\begin{frame}{Programaci\'on din\'amica: exactitud y resultados}
  \begin{columns}[T,onlytextwidth]
    \column{0.46\textwidth}
    \scriptsize
    \begin{block}{Resultados en los ejemplos}
      \begin{tabular}{@{}lcc@{}}
        \toprule
        Finca & \(\Pi^\star\) & Costo \\
        \midrule
        \(F_1\) & \(\langle 2,1,0,3,4\rangle\) & 20 \\
        \(F_2\) & \(\langle 2,1,0,3,4\rangle\) & 32 \\
        \bottomrule
      \end{tabular}
    \end{block}

    \begin{block}{Correcci\'on}
      La PD es exacta porque:
      \begin{itemize}
        \item el problema tiene subestructura \'optima;
        \item la recurrencia eval\'ua todas las transiciones posibles;
        \item la memoizaci\'on evita recomputar estados.
      \end{itemize}
    \end{block}

    \column{0.48\textwidth}
    \scriptsize
    \begin{block}{Evidencia experimental}
      En los 18 casos oficiales con salida de referencia:
      \begin{itemize}
        \item PD coincidi\'o con el costo \'optimo en todos los casos;
        \item esa columna se us\'o como referencia \'optima para medir el error del voraz.
      \end{itemize}
    \end{block}

    \begin{alertblock}{Utilidad pr\'actica}
      En el proyecto, PD fue viable hasta \(n\approx 20\).
      \begin{itemize}
        \item Para \(n=32\): \(137{,}438{,}953{,}472\) operaciones.
        \item Estimaci\'on pr\'actica: 458 min y 16 GB de memoria base.
      \end{itemize}
    \end{alertblock}
  \end{columns}
\end{frame}

\begin{frame}{Criterio voraz implementado}
  \begin{columns}[T,onlytextwidth]
    \column{0.42\textwidth}
    \scriptsize
    \begin{block}{Regla de orden}
      Regla base: \(\rho_i=\frac{tr_i}{p_i}\). Desempate: \(h_i=ts_i-tr_i\).
      \begin{enumerate}
        \setlength{\itemsep}{0.15em}
        \item Ordenar por \(\rho_i\) ascendente.
        \item Si empatan, ordenar por \(h_i\) ascendente.
        \item Si vuelven a empatar, conservar el menor \'indice.
      \end{enumerate}
      \centering
      \(\frac{tr}{p}\) balancea duraci\'on del riego y castigo por retraso.
      
      \medskip
      \textbf{Complejidad:} \(T_{\texttt{roV}}(n)=\Theta(n\log n)\).
    \end{block}

    \column{0.50\textwidth}
    \tiny
    \begin{block}{Implementaci\'on resumida}
      \begin{itemize}
        \item Construir \(orden=[0,1,\dots,n-1]\).
        \item Ordenar los \'indices con \texttt{sort}.
        \item Comparaci\'on principal por producto cruzado: \(tr_i p_j < tr_j p_i\).
        \item Si empatan, comparar holgura \(ts_i-tr_i\).
        \item Si vuelven a empatar, conservar el menor \'indice.
      \end{itemize}
    \end{block}
  \end{columns}
\end{frame}

\begin{frame}{L\'imite formal del voraz}
  \begin{columns}[T,onlytextwidth]
    \column{0.32\textwidth}
    \begin{block}{Resultado te\'orico}
      \begin{tabular}{@{}ll@{}}
        \toprule
        Propiedad & Estado \\
        \midrule
        Subestructura \'optima & S\'i \\
        Escogencia voraz & No \\
        \bottomrule
      \end{tabular}
    \end{block}

    \begin{alertblock}{Conclusi\'on}
      El problema admite soluciones exactas, pero el criterio voraz no garantiza costo \'optimo general.
    \end{alertblock}

    \column{0.60\textwidth}
    \small
    \begin{block}{Contraejemplo}
      \[
      T_0=\langle 10,2,1,0\rangle,
      \qquad
      T_1=\langle 5,3,4,0\rangle
      \]
      \[
      \frac{tr_0}{p_0}=2.00,
      \qquad
      \frac{tr_1}{p_1}=0.75
      \]
      \begin{tabular}{@{}lcc@{}}
        \toprule
        Orden & Regla voraz & Costo \\
        \midrule
        \(\langle 1,0\rangle\) & elegido por \(\frac{tr}{p}\) & 12 \\
        \(\langle 0,1\rangle\) & \'optimo real & 8 \\
        \bottomrule
      \end{tabular}
    \end{block}
  \end{columns}
\end{frame}

\begin{frame}{Bater\'ia del repositorio: 18 casos}
  \begin{columns}[T,onlytextwidth]
    \column{0.46\textwidth}
    \small
    \begin{block}{Resumen de resultados}
      \begin{tabular}{@{}lc@{}}
        \toprule
        M\'etrica & Valor \\
        \midrule
        Casos evaluados & 18 \\
        Coincidencias exactas en costo & \(1/18=5.56\%\) \\
        Diferencia promedio & 32.56 \\
        Exceso relativo promedio & 16.23\% \\
        \bottomrule
      \end{tabular}
    \end{block}

    \column{0.46\textwidth}
    \small
    \begin{block}{Peores casos}
      \begin{tabular}{@{}lcc@{}}
        \toprule
        Caso & Voraz & \'Optimo \\
        \midrule
        \texttt{test12} & 638 & 526 \\
        \texttt{test4} & 52 & 26 \\
        \bottomrule
      \end{tabular}
      \[
      d_k=\lvert C_V^{(k)}-C^{\star(k)}\rvert,
      \qquad
      e_k=100\frac{d_k}{C^{\star(k)}}
      \]
    \end{block}

    \begin{alertblock}{Lectura}
      La bater\'ia manual confirma fallas claras del voraz incluso en instancias peque\~nas.
    \end{alertblock}
  \end{columns}
\end{frame}

\begin{frame}{Validaci\'on emp\'irica: bloque mixto de 5000 fincas}
  \begin{columns}[T,onlytextwidth]
    \column{0.41\textwidth}
    \small
    \begin{block}{Dise\~no del experimento}
      \begin{tabular}{@{}lc@{}}
        \toprule
        Par\'ametro & Valor \\
        \midrule
        Tama\~no por finca & \(n=20\) \\
        Total de fincas & 5000 \\
        Total de tablones & 100000 \\
        Semilla & 20260513 \\
        \bottomrule
      \end{tabular}
    \end{block}

    \column{0.53\textwidth}
    \small
    \begin{block}{Resultado del voraz}
      \begin{tabular}{@{}lc@{}}
        \toprule
        M\'etrica & Valor \\
        \midrule
        Coincidencias exactas en costo & \(219/5000=4.38\%\) \\
        Diferencia absoluta promedio & 66.94 \\
        Exceso relativo promedio & 1.97\% \\
        Peor diferencia observada & 382 \\
        \bottomrule
      \end{tabular}
    \end{block}

    \begin{alertblock}{Lectura}
      En mezcla realista, el voraz rara vez iguala el costo \'optimo, pero su exceso relativo promedio sigue siendo bajo.
    \end{alertblock}
  \end{columns}
\end{frame}

\begin{frame}{Experimento de aislamiento por perfiles puros}
  \scriptsize
  \begin{table}
    \centering
    \begin{tabularx}{\textwidth}{@{}>{\RaggedRight\arraybackslash}X>{\Centering\arraybackslash}p{2.2cm}>{\Centering\arraybackslash}p{2.3cm}>{\Centering\arraybackslash}p{2.7cm}@{}}
      \toprule
      Perfil exclusivo & Dif. promedio & Exceso relativo & Coincidencia exacta \\
      \midrule
      \(100\%\) Cr\'iticos & 0.00 & 0.00\% & \(5000/5000\) \\
      \(100\%\) Urgentes & 2.83 & 0.05\% & \(1735/5000\) \\
      \(100\%\) Balanceados & 11.13 & 0.35\% & \(533/5000\) \\
      \(100\%\) Largo plazo & 180.25 & 1.41\% & \(25/5000\) \\
      \(100\%\) Relajados & 62.37 & 5.53\% & \(4/5000\) \\
      \bottomrule
    \end{tabularx}
  \end{table}

  \vspace{0.3em}
  \begin{columns}[T,onlytextwidth]
    \column{0.47\textwidth}
    \begin{block}{Acierta mejor}
      Cr\'itico y urgente: cuando retrasarse casi siempre empuja al caso 3.
    \end{block}

    \column{0.47\textwidth}
    \begin{alertblock}{Falla m\'as}
      Relajado y largo plazo: cuando la holgura y el d\'ia perfecto pesan m\'as que la urgencia.
    \end{alertblock}
  \end{columns}
\end{frame}

\begin{frame}{Comparaci\'on de criterios voraces}
  \scriptsize
  \begin{table}
    \centering
    \begin{tabular}{@{}l l c c@{}}
      \toprule
      Principal & Desempate & \% dif. & Exactas \\
      \midrule
      \(\frac{tr}{p}\) & \(ts-tr\) & 1.97\% & 219/5000 \\
      \(\frac{tr}{p}\) & supervivencia & 1.98\% & 215/5000 \\
      \(\frac{tr}{p}\) & d\'ia perfecto & 2.08\% & 168/5000 \\
      \(\frac{tr}{p}\) & prioridad alta & 2.20\% & 143/5000 \\
      supervivencia & \(\frac{tr}{p}\) & 26.14\% & 0/5000 \\
      \bottomrule
    \end{tabular}
  \end{table}

  \vspace{0.2em}
  \begin{block}{Lectura}
    Las cuatro mejores variantes mantienen \(\frac{tr}{p}\) como criterio principal.
  \end{block}

  \begin{alertblock}{Conclusi\'on}
    El desempate importa, pero la decisi\'on clave es conservar \(\frac{tr}{p}\) como criterio base. El mejor desempate fue \(ts-tr\).
  \end{alertblock}
\end{frame}

\begin{frame}{Conclusiones}
  \small
  \begin{block}{Mensaje central}
    No existe una estrategia universalmente mejor: cada enfoque responde a una necesidad distinta de exactitud y costo computacional.
  \end{block}

  \begin{enumerate}
    \item \textbf{Fuerza bruta:} garantiza el \'optimo y sirve como referencia exacta, pero su costo factorial la limita a \(n\) peque\~no.
    \item \textbf{Programaci\'on din\'amica:} tambi\'en es exacta y mejora dr\'asticamente la escalabilidad hasta instancias moderadas.
    \item \textbf{Voraz:} \(\Theta(n\log n)\), muy r\'apido, pero sin garant\'ia de costo \'optimo general.
    \item \textbf{Teor\'ia + evidencia:} el voraz acierta mejor cuando domina el riesgo de tardanza y falla m\'as cuando la holgura y \(rp\) pesan m\'as.
    \item \textbf{Selecci\'on final:} entre las variantes probadas, \(\frac{tr}{p}\) con desempate por \(ts-tr\) fue la m\'as robusta.
    \item \textbf{Uso recomendado:} si se necesita exactitud, usar PD; si se necesita rapidez inmediata, usar el voraz.
  \end{enumerate}
\end{frame}

\end{document}
